\section{Certified descent to the conjugation-fixed fields}
\label{sec:cubic-real-descent}

Write $\alpha^3=269$, $F=\mathbb Q(\alpha)$, and
$F_1=\mathbb Q(\alpha,\zeta_9+\zeta_9^{-1})$.
Here $\alpha$ denotes the real root in the chosen complex embedding.
The fields $F$ and $F_1$ have signatures $(1,1)$ and $(3,3)$;
``conjugation-fixed'' does not mean that all their embeddings are real.
The involution $c$ is the actual complex conjugation in
$E=F(\omega)$ and $E_1=F_1(\omega)$.

\begin{proposition}[AC-K22: certified real-field descent and a nonconstant pattern]
\label{AC-K22}
There are natural isomorphisms
\[
 \widetilde{\operatorname{Cl}}_3(F)
       \simeq\widetilde{\operatorname{Cl}}_3(E)^{c=1},\qquad
 \widetilde{\operatorname{Cl}}_3(F_1)
       \simeq\widetilde{\operatorname{Cl}}_3(E_1)^{c=1}.
                                                        \tag{KD.1}
\]
The second isomorphism respects the cyclotomic action
$\gamma(\zeta_9)=\zeta_9^4$, $\gamma(\alpha)=\alpha$.
Unconditionally certified calculations give
\[
 \widetilde{\operatorname{Cl}}_3(F)=\mathbb Z/3,
 \qquad
 \widetilde{\operatorname{Cl}}_3(F_1)=\mathbb Z/9.
                                                        \tag{KD.2}
\]
Consequently the all-four alternative in AC-K21 does not occur:
there is a fixed $r\in\{1,2\}$ such that
\[
                m_e=4+1_{e\equiv r\pmod3}
                   \qquad(e\geq1).                    \tag{KD.3}
\]
The certification of (KD.2) alone does not decide whether $r=1$
or $r=2$.  The remaining datum is the oriented action of
$\gamma$ on the cyclic group in the second part of (KD.2).
If that action is multiplication by $1+3c_0$ modulo nine,
then $c_0\in\{1,2\}$ and $r\equiv-c_0\pmod3$.
\end{proposition}
\begin{proof}
\proofstep{1. Prime-to-three descent preserves degree zero.}
For a quadratic Galois extension $L/K$ with involution $c$,
let $j$ and $N$ denote extension and norm on logarithmic
three-adic divisor classes.  Use ordinary local ramification and inertia degrees away
from three, and logarithmic local degrees at three.
This is the computational normalization in Belabas--Jaulent,
\S1, pp.~2--3: its logarithmic valuations are the ordinary
valuations away from three.  Denote the chosen local pair
by $e_*(w/v),f_*(w/v)$; in both cases
$e_*(w/v)f_*(w/v)=[L_w:K_v]$.
On a prime divisor, extension with coefficients $e_*$
and norm with coefficients $f_*$ therefore give
\[
                      Nj=2,\qquad jN=1+c.             \tag{KD.4}
\]
The same formulas hold on principal logarithmic divisors:
they express compatibility of the local logarithmic
valuation with field norm, and with the product of the
conjugates.  Extension and norm preserve the degree-zero
submodules, since logarithmic degree is compatible with
these local degrees.  Thus (KD.4) holds on the degree-zero
logarithmic class groups, not just on their full class groups.
These are the logarithmic divisor and norm conventions of
AC-K13 and its Belabas--Jaulent source; this argument uses
only their defining functoriality.

Multiplication by two is invertible on a three-primary
$\mathbb Z_3$-module.  The first identity makes $j$ injective.
If $z$ is fixed by $c$, the second identity gives
$z=j(\tfrac12 Nz)$.  The image of $j$ is fixed by $c$.
This proves the asserted descent for any such quadratic
extension, without a conjecture about finiteness or an
assumption of logarithmic unramifiedness.
Apply it to $E/F$ and $E_1/F_1$.
All the defining maps commute with automorphisms that
commute with $c$, in particular the stated $\gamma$;
this proves the equivariance in (KD.1).

\proofstep{2. The exact arithmetic certificates.}
For $F$, the defining polynomial is $y^3-269$.
PARI/GP~2.19.1 returned signature $[1,1]$, discriminant
$-217083$, and ordinary class group $[7]$.
The full command \texttt{bnfcertify(b)} returned \texttt{1}.
The command \texttt{bnflog(b,3)} then returned
\texttt{[[3],[3],[]]}.
For $F_1$, the script forms
\texttt{polcompositum(y\string^3-269,y\string^3-3*y+1)[1]}
and applies \texttt{polredbest}; its resulting polynomial is
\[
 \begin{split}
 p(y)={}&y^9-3y^7-270y^6+810y^5+1347y^4\\
        &-48600y^3-72900y^2-96123y-728731.
 \end{split}                                         \tag{KD.5}
\]
Its signature is $[3,3]$, discriminant
$-604073788222005661563$, and ordinary class group $[21,3]$.
Again the full \texttt{bnfcertify(b)} returned \texttt{1},
and \texttt{bnflog(b,3)} returned \texttt{[[9],[3],[3]]}.
In each case the first component is the degree-zero
logarithmic class group.  The certification and the
logarithmic algorithm have exactly the meaning fixed in
AC-K13: Belabas--Jaulent, \emph{The logarithmic class group
package in PARI/GP}, author version \S\S3.2--3.4,
Lemma~3.8 and \S4.1,
\url{https://www.math.u-bordeaux.fr/~kbelabas/research/bnflog-2016.pdf}.
No partial-certification flag is used.

For an additional exact field-identification check,
\texttt{polisirreducible(p)} returns one and
\texttt{nfcertify(b)} returns the empty vector.
The saved script produces elements $a,b_0$ of
$\mathbb Q[y]/(p)$ and verifies exactly
$a^3=269$ and $b_0^3-3b_0+1=0$.
The cyclic cubic field defined by the second equation
is $\mathbb Q(\zeta_9+\zeta_9^{-1})$.
It is distinct from the non-Galois pure cubic field,
so their compositum has degree nine.
Thus these inclusions identify the computational field
with $F_1$, without relying only on a numerical embedding.
The same script finds an exact field automorphism fixing
$a$ and sending $b_0$ to $2-b_0-b_0^2$.
For $b_0=\zeta_9+\zeta_9^{-1}$ this is precisely
$\zeta_9^4+\zeta_9^{-4}$; hence it specifies the orientation
needed for a later action certificate.
It does not yet compute that action on the logarithmic group.

The scripts and complete raw outputs are
\path{evidence/cubic_logarithmic/pure_cubic_base.gp},
\path{pure_cubic_base.out},
\path{real_cubic_layer1.gp}, and
\path{real_cubic_layer1.out}.
The last script's first logarithmic output was labelled
preliminary; the later full certification and repeated
logarithmic output establish (KD.2).
The additional exact field checks are retained separately
as \path{real_cubic_field_action.gp} and its output.
The runtime build receipt and algorithm references are
unchanged from AC-K13 and AC-K15.

\proofstep{3. Identify the small component and apply AC-K21.}
By AC-K13, the base group for $E$ is
$\mathbb Z/9\oplus\mathbb Z/3$.
Equations (KD.1)--(KD.2) show that its actual $c$-positive
component has order three.  Its order-nine component
therefore has negative $c$ sign.
Accordingly the generator $x$ in AC-K18--AC-K21 lies
in the actual negative-sign component; the opposite
component $U$ used there is the actual positive-sign component.
The polynomial $F(B)=\gamma^3-1$ annihilates $U$ in every
case proved there.  Schneider's logarithmic descent and
(KD.1) consequently give the exact identifications
\[
 U=(M/F(B)M)^{c=1}
   =\widetilde{\operatorname{Cl}}_3(E_1)^{c=1}
   \simeq\widetilde{\operatorname{Cl}}_3(F_1)
   \simeq\mathbb Z/9.                                \tag{KD.6}
\]
Thus the elementary-component branch of AC-K21 is excluded.
Its cyclic-nine branch proves (KD.3) and the stated
oriented-action rule.  The Tate twist may interchange
conjugation labels on a specialized module; the comparison
here uses the original untwisted $M$ throughout, as in AC-K21.
\end{proof}

The exact fixed-three contribution for $N=t-1$ is therefore
\[
 \sum_{e=1}^{N}m_e\log3
   =\left(4N+\left\lfloor\frac{N+3-r}{3}\right\rfloor\right)\log3,
                      \qquad r\in\{1,2\}.            \tag{KD.7}
\]
The subsequent labelled-action certificate AC-K23 proves
$r=2$ and thereby resolves the remaining choice in (KD.7).
This selects the nonconstant alternative using certified
arithmetic; it does not supply the other-prime aggregate
or identify an automorphic theta object with the source Selmer group.
The norm/extension argument is a finite Galois descent
calculation.  In the finite-\'etale dictionary already used
from SGA~1, Expos\'e~V, \S\S4 and~7, it is the invariant
summand for an involution of order invertible on the coefficients;
no additional geometric realization or archive locator is asserted.
